\subsubsection{Coordinate Transformations and Sobolev Spaces on Manifolds}\label{subsec:sobolev-coordinate-transformations-manifolds}

If $\partial\Omega$ is smooth, a standard trick is to make local changes of coordinates which make $\partial\Omega$ a coordinate surface. This requires us to consider how Sobolev spaces behave under coordinate transformations. Also, for boundary-value problems in PDEs we need to consider spaces of functions defined on $\partial\Omega$. Hence we want to define Sobolev spaces on manifolds. These will of course be defined in terms of local coordinate charts; again the behavior under coordinate transformations is a crucial issue.

\begin{Definition}
\label{def:sobolev-k-diffeomorphism}
Let $\Omega,\Omega^{\prime}$ be open sets in $\mathbb{R}^m$. A bijection $\bm{\Phi}:\Omega\to\Omega^{\prime}$ is called a $k$-diffeomorphism $(k\geq1)$ if\begin{enumerate}
\item $\bm{\Phi}$ and $\bm{\Phi}^{-1}$ are continuous on $\overline\Omega$ and $\overline{\Omega^{\prime}}$, respectively,
\item their derivatives of order $1$ through $k$ are bounded and continuous on $\overline\Omega$ and $\overline{\Omega^{\prime}}$, respectively, and
\item there are positive constants $c$ and $C$ such that $c\leq|\det\nabla\bm{\Phi}|\leq C$ in $\Omega$.
\end{enumerate}
\end{Definition}

\begin{Definition}
\label{def:sobolev-diffeomorphism-pullbacks}
Let $\bm{\Phi}:\Omega\to\Omega^{\prime}$ be a $k$-diffeomorphism. Then the pullback operators $\bm{\Phi}^*$ and $(\bm{\Phi}^{-1})^*$ are defined by
\begin{equation}\label{eq:sobolev-diffeomorphism-pullback-definitions}
\bm{\Phi}^*u=u\circ\bm{\Phi},\quad(\bm{\Phi}^{-1})^*v=v\circ\bm{\Phi}^{-1}
\end{equation}
for functions $u,v$ defined, respectively, on $\Omega^{\prime}$ and $\Omega$.
\end{Definition}

\begin{Theorem}
\label{thm:sobolev-integer-order-pullback}
Let $\bm{\Phi}:\Omega\to\Omega^{\prime}$ be a $k$-diffeomorphism. Then $\bm{\Phi}^*$ and $(\bm{\Phi}^{-1})^*$ are bounded linear mappings from $H^k(\Omega^{\prime})$ to $H^k(\Omega)$ and, respectively, from $H^k(\Omega)$ to $H^k(\Omega^{\prime})$.
\end{Theorem}

\begingroup
\renewcommand{\endProof}{\par}
\begin{Proof}
Since the inverse of a $k$-diffeomorphism is also a $k$-diffeomorphism, it suffices to prove the claim for $\bm{\Phi}^*$. It is easy to see that $\bm{\Phi}^*$ is a bounded linear mapping from $\widetilde L^2(\Omega^{\prime})$ to $\widetilde L^2(\Omega)$ and hence, by continuity, from $L^2(\Omega^{\prime})$ to $L^2(\Omega)$. We note that $(u,v)_{\Omega^{\prime}}=(\bm{\Phi}^*u,(\det\nabla\bm{\Phi})\bm{\Phi}^*v)_{\Omega}$.

It remains to consider derivatives. For this, we have to establish that the derivatives of $\bm{\Phi}^*u$ can be evaluated in the usual way by means of the chain and product rule. This is clear if $u\in C^{\infty}(\Omega)$, and we have already established that $C^{\infty}$-functions are dense in $H^k(\Omega)$. For $u\in H^k(\Omega)\cap C^{\infty}(\Omega)$, any $p$th derivative of $\bm{\Phi}^*u$ $(p\leq k)$ is a linear expression involving derivatives of $u$ of orders $\leq p$, with coefficients depending on derivatives of $\bm{\Phi}$; in particular, all the coefficients are bounded, continuous functions. In order to extend the result to all of $H^k(\Omega)$, we need to be able to take limits. The following lemma establishes that we can do this.
\end{Proof}
\endgroup

\begin{Lemma}
\label{lem:l-two-bounded-continuous-multiplier}
Let $u\in L^2(\Omega)\cap C(\Omega)$ and $a\in C_b(\Omega)$. Then $au\in L^2(\Omega)$ and $\|au\|_2\leq\|a\|_{\infty}\|u\|_2$.
\end{Lemma}

\begingroup
\postdisplaypenalty=10000
\renewcommand{\endProof}{\par\nobreak\hfill$\blacksquare$\par\vspace{0.1cm}}
\begingroup
\renewenvironment{Proof}{\par}{\par\nobreak\hfill$\blacksquare$\par\vspace{0.1cm}}
\begin{Proof}
We omit the trivial proof of the lemma. Since $L^2(\Omega)\cap C(\Omega)$ is dense in $L^2(\Omega)$, the lemma allows us to define $au$ for every $u\in L^2(\Omega)$. This completes the proof of Theorem~\ref{thm:sobolev-integer-order-pullback}.
\end{Proof}
\endgroup
\endgroup
